\section{Background}
\label{sec:background}

\subsection{Notation}
Let a sequence $x=(x_1,\ldots,x_L)$ be divided into $B$ contiguous blocks of
$M=L/B$ tokens.
\begin{equation}
    x=\left(x^{(1)},\ldots,x^{(B)}\right),
    \qquad
    x^{(<b)}=\left(x^{(1)},\ldots,x^{(b-1)}\right).
\end{equation}
We use $b$ to index blocks and $r$ to index the $P$ ranks in a parallel group.
At each transformer layer~\citep{vaswani2017attention}, $Q$, $K$, and $V$ denote the query, key, and value
projections; semicolons inside brackets denote sequence concatenation.

\subsection{Block diffusion language models}

A block diffusion language model (BDLM) generates blocks from left to right
while denoising the tokens within each block in parallel
\citep{arriola2025blockdiffusion}.
\begin{equation}
    p_\theta(x)
    =
    \prod_{b=1}^{B}
    p_\theta\!\left(x^{(b)}\mid x^{(<b)}\right),
    \label{eq:bdlm-factorization}
\end{equation}
where each conditional distribution is modeled by a diffusion process.

% Verified against bdlm_parallel commit 4b60e48: core/objectives/runtime.py
% (stratified times), core/objectives/fast_dllm_v2.py (uniform times), and
% training/run_spec.py (default sampling bounds).
Training samples continuous times $t_b$, not discrete timestep indices:
standard BDLM uses uniform stratification over $[10^{-3},1)$ across batch/block
pairs; Fast-dLLM~v2 uses unstratified $\mathcal{U}[0,1)$ draws.
Each target block forms a corrupted copy $x_{t_b}^{(b)}$ conditioned on its
clean prefix $x^{(<b)}$. The training objective is
\begin{equation}
    \mathcal{L}_{\mathrm{BDLM}}
    =
    \sum_{b=1}^{B}
        \mathbb{E}_{t_b,\,x_{t_b}^{(b)}}\!\left[
            -w(t_b)\log p_\theta\!\left(
                x^{(b)} \mid x_{t_b}^{(b)},x^{(<b)}
            \right)
        \right].
    \label{eq:bdlm-training-objective}
\end{equation}
The expectation is over $t_b$ and
$x_{t_b}^{(b)}\sim q_{t_b}(\cdot\mid x^{(b)})$, where $q_{t_b}$ is the
corruption process and $w(t_b)$ is a noise-dependent loss weight. We denote
the $b$th summand, including token normalization, by $\mathcal{L}_b$. Training
observes the complete sequence and uses the clean ground-truth prefix for every
target block.

The standard BDLM training computation evaluates all $B$ block losses together
by placing a length-$L$ corrupted copy alongside a length-$L$ clean copy,
forming a length-$2L$ representation. It contains a block-specific corrupted
copy and a reusable clean copy of every token. A block-causal attention mask
gives the clean tokens in block $b$ access to the preceding clean blocks and
causal access within block $b$. The corrupted tokens in target block $b$
attend bidirectionally within their own corrupted block and attend to the clean
prefix $x^{(<b)}$, but do not use corrupted K/V from other blocks. Their
visible keys and values are therefore
$[K_b;K_{<b}^{\mathrm{clean}}]$ and
$[V_b;V_{<b}^{\mathrm{clean}}]$.

The corrupted activations of block $b$ affect only $\mathcal{L}_b$, while its
clean representation can be reused by multiple later block losses. Corrupted
computation is therefore block-specific, while clean computation is shared
across block losses. All block losses share the model parameters.


\subsection{Context parallelism}
\label{sec:background-cp}
Context parallelism (CP) uses token position as its ownership dimension,
partitioning the attention input across $P$ ranks so each rank stores
approximately a $1/P$ fraction of its token activations while the group jointly
evaluates attention~\citep{liu2023ringattention,jacobs2023ulysses}. BDLM
training supplies the length-$2L$ clean-plus-corrupted representation as this
attention input, so conventional CP partitions both copies by token position.

Let $Q_r$, $K_r$, and $V_r$ denote the query, key, and value projections
initially owned by rank $r$, and let $d_h$ be the attention-head dimension.
Writing $K=[K_1;\ldots;K_P]$ and $V=[V_1;\ldots;V_P]$, the attention output
for the queries on rank $r$ is
\begin{equation}
O_r
=
\operatorname{softmax}\!\left(
    \frac{Q_rK^\top}{\sqrt{d_h}} + \mathcal{M}_r
\right)V,
\label{eq:context-parallel-attention}
\end{equation}
where $\mathcal{M}_r$ contains the attention-mask rows for the queries on rank
$r$.

Causal attention produces uneven work across sequence shards because later
queries see longer prefixes. Zigzag partitioning pairs early and late chunks
on each rank to balance their attention work~\citep{dubey2024llama3}, while
retaining the sequence shard as the unit of CP ownership.
